% ============================================================
% 设备完整性在线监测系统 — 网络拓扑与访问权限白名单
% ============================================================
\documentclass[12pt,a4paper]{ctexart}

% ── 编码与字体 ──
\usepackage{fontspec}
\usepackage{xeCJK}
\setCJKmainfont{Microsoft YaHei}
\setmainfont{Calibri}

% ── 页面 ──
\usepackage{geometry}
\geometry{top=2.5cm, bottom=2.5cm, left=2.8cm, right=2.8cm}

% ── 常用包 ──
\usepackage{graphicx}
\usepackage{float}
\usepackage{fancyhdr}
\usepackage{hyperref}
\usepackage{longtable}
\usepackage{booktabs}
\usepackage{array}
\usepackage{xcolor}
\usepackage{caption}
\usepackage{enumitem}

% ── 颜色定义 ──
\definecolor{green}{HTML}{2E7D32}
\definecolor{orange}{HTML}{E65100}
\definecolor{red}{HTML}{C62828}
\definecolor{blue}{HTML}{1565C0}
\definecolor{grey}{HTML}{616161}

% ── 状态标签 ──
\newcommand{\statusok}{\textcolor{green}{\textbf{已通}}}
\newcommand{\statuspending}{\textcolor{orange}{\textbf{待申请}}}
\newcommand{\statusapprove}{\textcolor{orange}{\textbf{待审批}}}

% ── 页眉页脚 ──
\pagestyle{fancy}
\fancyhead[L]{设备完整性在线监测系统}
\fancyhead[R]{网络拓扑与访问权限白名单}
\fancyfoot[C]{\thepage}

% ── 文档信息 ──
\title{\textbf{设备完整性在线监测系统}\\网络拓扑与访问权限白名单}
\author{迪格(杭州)物联科技有限公司}
\date{V1.0 \quad 2026-07-17}

\begin{document}
\maketitle
\thispagestyle{fancy}

% ════════════════════════════════════════════════════════════
\section{网络拓扑总览}
% ════════════════════════════════════════════════════════════

系统网络分为三个安全区域：\textbf{场站内网}（数据源头）、\textbf{生产网/分厂}（数据汇聚与展示）、\textbf{DMZ}（远程访问），通过防火墙实现安全隔离。

\begin{figure}[H]\centering
\small
\begin{tabular}{p{0.28\textwidth} p{0.04\textwidth} p{0.28\textwidth} p{0.04\textwidth} p{0.28\textwidth}}
\toprule
\centering\textbf{场站内网} \newline 192.168.x.x & & \centering\textbf{生产网/分厂} \newline 10.99.0.0/24 & & \centering\textbf{DMZ} \newline 11.66.113.x/25 \\
\midrule
\centering\footnotesize
传感器/PLC/NVR \newline
边缘采集服务器 \newline
海康 NVR

\newline\newline
\textbf{数据源头} & $\longrightarrow$ &
\centering\footnotesize
DG-IoT 采集服务器 \newline
render 01{\raise.17ex\hbox{$\sim$}}05 \newline
渲染+采集+大屏

\newline\newline
\textbf{数据汇聚+展示} & $\longrightarrow$ &
\centering\footnotesize
大屏/远程访问 \newline
:5173 :5050 \newline
内部 IP 访问

\newline\newline
\textbf{对外服务} \\
\bottomrule
\multicolumn{5}{c}{\footnotesize 端口: 502/1883/554/443}
\end{tabular}
\caption{三区网络拓扑}
\end{figure}

\subsection{数据流向}

\begin{enumerate}[leftmargin=*]
  \item \textbf{生产网采集服务器} $\longrightarrow$ 防火墙 $\longrightarrow$ \textbf{场站内网}
  \begin{itemize}
    \item Modbus TCP :502 $\longrightarrow$ PLC（锅炉能效）
    \item RTSP :554 $\longrightarrow$ 海康 NVR（视频/TDLAS）
    \item MQTT :1883 $\longrightarrow$ 博感网关（智能螺栓）
    \item HTTPS :443 $\longrightarrow$ 边缘服务器 API（油液/振温）
  \end{itemize}
  \item \textbf{分厂 (10.99.0.x)} $\longrightarrow$ 白名单 :443 $\longrightarrow$ \textbf{DMZ (11.66.113.x)}
  \begin{itemize}
    \item 领导在办公网浏览器打开 \texttt{http://11.66.113.78:5173}
  \end{itemize}
\end{enumerate}

% ════════════════════════════════════════════════════════════
\section{防火墙白名单}
% ════════════════════════════════════════════════════════════

\subsection{生产网 $\rightarrow$ 场站方向}

\begin{longtable}{p{0.7cm} p{2.2cm} p{1.2cm} p{0.8cm} p{2.5cm} p{2.5cm} p{1.5cm}}
\toprule
\textbf{序号} & \textbf{应用} & \textbf{协议} & \textbf{端口} & \textbf{源地址} & \textbf{目标地址} & \textbf{状态} \\
\midrule
\endfirsthead
\multicolumn{7}{c}{\small（续表）} \\
\toprule
\textbf{序号} & \textbf{应用} & \textbf{协议} & \textbf{端口} & \textbf{源地址} & \textbf{目标地址} & \textbf{状态} \\
\midrule
\endhead
\bottomrule
\endfoot
1 & 锅炉能效（跃佳 YJ-NX-1） & Modbus TCP & 502 & 10.99.0.x/24 & 192.168.x.x & \statuspending \\
2 & 视频流（海康 NVR） & RTSP & 554 & 10.99.0.x/24 & 192.168.x.x & \statuspending \\
3 & 螺栓监测（博感 GU100X） & MQTT & 1883 & 10.99.0.x/24 & 192.168.x.x & \statusok \\
4 & 边缘API/有叶云 & HTTPS & 443 & 10.99.0.x/24 & 192.168.x.x & \statusok \\
5 & WinRM 管理（主站） & HTTP & 5985 & 11.66.12.131 & 11.66.12.131 & 内网直连 \\
\end{longtable}

\subsection{分厂 $\rightarrow$ DMZ 方向}

\begin{longtable}{p{0.7cm} p{2.5cm} p{0.8cm} p{2.5cm} p{2.5cm} p{1.5cm}}
\toprule
\textbf{序号} & \textbf{服务} & \textbf{端口} & \textbf{源地址} & \textbf{目标地址} & \textbf{状态} \\
\midrule
\endfirsthead
\multicolumn{6}{c}{\small（续表）} \\
\toprule
\textbf{序号} & \textbf{服务} & \textbf{端口} & \textbf{源地址} & \textbf{目标地址} & \textbf{状态} \\
\midrule
\endhead
\bottomrule
\endfoot
1 & 前端大屏 & 5173 & 10.99.0.x/24 & 11.66.113.x/25 & \statusapprove \\
2 & 后端 API & 5050 & 10.99.0.x/24 & 11.66.113.x/25 & \statusapprove \\
3 & EMQX Dashboard & 18083 & 10.99.0.x/24 & 11.66.113.x/25 & \statusapprove \\
\end{longtable}

% ════════════════════════════════════════════════════════════
\section{访问权限矩阵}
% ════════════════════════════════════════════════════════════

\begin{longtable}{p{0.5cm} p{1.8cm} p{3cm} p{3cm} p{3.5cm} p{2cm}}
\toprule
\textbf{序号} & \textbf{所属单位} & \textbf{站场} & \textbf{可访问设备} & \textbf{数据上送目标} & \textbf{视频流开放} \\
\midrule
\endfirsthead
\multicolumn{6}{c}{\small（续表）} \\
\toprule
\textbf{序号} & \textbf{所属单位} & \textbf{站场} & \textbf{可访问设备} & \textbf{数据上送目标} & \textbf{视频流开放} \\
\midrule
\endhead
\bottomrule
\endfoot
1 & 天然气 & 北1-2天然气处理厂 & 已接入的一次/二次设备、NVR、物联网网关全部采集数据 & — & — \\
2 & 天然气 & 北1-2 & 设备完整性采集数据 & $\rightarrow$ 天然气分公司信息中心 & — \\
3 & 天然气 & 北1-2天然气处理厂 & — & — & 天然气分公司信息中心可访问新装摄像头视频流 \\
4 & 采油二厂 & 南4联合站 & 已接入的一次/二次设备、NVR、物联网网关全部采集数据 & — & — \\
5 & 采油二厂 & 南4联合站 & 设备完整性采集数据 & $\rightarrow$ 采油二厂第四作业区信息中心 & — \\
6 & 采油二厂 & 南4联合站 & — & — & 第四作业区信息中心可访问新装摄像头视频流 \\
7 & 采油三厂 & 北9注水站/北15联合站/萨北21站 & 各站已接入的一次/二次设备、NVR、物联网网关全部采集数据 & — & — \\
8 & 采油三厂 & 北9/北15/萨北21 & 设备完整性采集数据 & $\rightarrow$ 采油三厂第八作业区信息中心 & — \\
9 & 采油三厂 & 北9/北15/萨北21 & — & — & 第八作业区信息中心可访问新装摄像头视频流 \\
\end{longtable}

% ════════════════════════════════════════════════════════════
\section{关键 IP 地址}
% ════════════════════════════════════════════════════════════

\begin{longtable}{p{4cm} p{3cm} p{6cm}}
\toprule
\textbf{系统} & \textbf{IP 地址} & \textbf{用途} \\
\midrule
\endhead
主站（pSpace） & 11.66.12.131 & WinRM 管理 :5985 \\
IO 服务器（A11） & 11.66.12.130 & A11 协议 :8889 \\
生产网段 & 10.99.0.0/24 & 采集服务器/渲染工作站 \\
DMZ 网段 & 11.66.113.x/25 & 大屏/远程访问（参考 11.66.113.78） \\
场站内网 & 192.168.x.x & PLC/NVR/网关等现场设备 \\
\end{longtable}

% ════════════════════════════════════════════════════════════
\section{端口汇总}
% ════════════════════════════════════════════════════════════

\begin{longtable}{p{1cm} p{1.5cm} p{3.5cm} p{3.5cm} p{2cm}}
\toprule
\textbf{端口} & \textbf{协议} & \textbf{用途} & \textbf{方向} & \textbf{状态} \\
\midrule
\endhead
\bottomrule
\endfoot
502 & Modbus TCP & 锅炉能效（YJ-NX-1） & 生产网$\rightarrow$场站 & \statuspending \\
554 & RTSP & 视频流（海康 NVR） & 生产网$\rightarrow$场站 & \statuspending \\
1883 & MQTT & 螺栓监测（博感 GU100X） & 生产网$\rightarrow$场站 & \statusok \\
443 & HTTPS & 边缘API/有叶云 & 生产网$\rightarrow$场站 & \statusok \\
5985 & HTTP & WinRM 管理 & 内网 & 内网直连 \\
5173 & HTTP & 前端大屏 & 分厂$\rightarrow$DMZ & \statusapprove \\
5050 & HTTP & 后端API & 分厂$\rightarrow$DMZ & \statusapprove \\
18083 & HTTP & EMQX 管理 & 分厂$\rightarrow$DMZ & \statusapprove \\
\end{longtable}

% ════════════════════════════════════════════════════════════
\section{摄像头接入方案}
% ════════════════════════════════════════════════════════════

\subsection{数据流向}

\begin{enumerate}[leftmargin=*]
  \item 海康 NVR 位于场站内网（192.168.x.x），摄像头通过局域网接入 NVR
  \item 生产网采集服务器通过 RTSP 协议（:554）从 NVR 拉取视频流，需防火墙开 554 端口
  \item 每个场站部署边缘工作站，运行 DG-IoT RTSP 代理服务
  \item DMZ 大屏通过 HLS/WebSocket 转推方式查看实时视频
\end{enumerate}

\subsection{RTSP 流地址格式}

\begin{verbatim}
rtsp://[user:pass@][NVR_IP]:554/Streaming/Channels/[channel]01
\end{verbatim}

其中 channel 编号规则：101=第1路主码流，102=第1路子码流，201=第2路主码流，以此类推。

\subsection{摄像头清单}

| 型号 | 类型 | 数量 | 安装场站 |
|------|------|:----:|----------|
| iDS-2DF84ADI-CWX | 防爆球机 | 11 | 北1-2$\times$9、北9$\times$1、北15$\times$1 |
| HM-TD25C8T-B | 热成像枪机 | 5 | 北1-2$\times$3、北15$\times$1、萨北21$\times$1 |
| HM-TD66AC-D | 双光谱云台 | 3 | 北1-2$\times$2、北9$\times$1 |
| DS-2XE94SH-LY | TDLAS 防爆 | 4 | 北1-2$\times$1、北15$\times$1、萨北21$\times$1 |

% ════════════════════════════════════════════════════════════
\section{版本记录}
% ════════════════════════════════════════════════════════════

\begin{table}[H]\centering
\begin{tabular}{p{1.5cm} p{2cm} p{2cm} p{7cm}}
\toprule
\textbf{版本} & \textbf{日期} & \textbf{修订人} & \textbf{修订内容} \\
\midrule
V1.0 & 2026-07-17 & — & 初始版本，基于现场调研和防火墙规则确认 \\
\bottomrule
\end{tabular}
\end{table}

\end{document}
